\section{Related Work}
\label{sec:related}

\paragraph{Multi-embodiment pretraining.}
Pooled demonstration corpora across institutions and robot types \citep{oxe2023openx, khazatsky2024droid} enabled generalist policies that adapt to new platforms after task-specific data collection \citep{kim2024openvla, ghosh2024octo, black2024pi0, pi2025pi05}. Subsequent systems widened the range of embodiments to navigation and locomotion \citep{doshi2024crossformer, yang2024pushing} and to human-centric sources \citep{luo2026beingh05}. These works vary data scale, embodiment coverage and architecture together, so the contribution of any single factor to transfer stays entangled with the rest of the system. Our study fixes the backbone, the budget and the evaluation and varies one channel at a time.

\paragraph{Cross-embodiment policy design.}
Policies that transfer across robots share either the interface or the supervision. Interface-level approaches learn latent or tokenised action spaces \citep{mu2026onepolicy, zheng2025xvla} or mask the end effector from the image \citep{piseno2026cloak}, and handheld collection devices remove the robot from data collection altogether \citep{chi2024umi, liu2026rdt2}. Supervision-level approaches add language descriptions of motion \citep{zha2026lap} or drop proprioceptive input \citep{zhao2025proprio}. Camera-centric action geometry, in which the policy predicts the motion of the end effector in image coordinates and an embodiment-specific translator produces commands, was introduced for large-scale heterogeneous pretraining \citep{xiaomi2026ucagp} and is related to keypoint-based action targets \citep{xu2026actiongeometry}. Our effect representation shares the premise that the camera frame is the common ground of embodiments and differs in three respects: the origin of every keypoint is subtracted analytically, the decoder is a closed-form geometric solve in place of a learned translator, and the representation is compared with base-frame and tool-frame alternatives under a controlled protocol that separates it from data composition.

\paragraph{Benchmarks and protocols.}
Cross-embodiment evaluations in the literature mix the presence of the target robot in pretraining with its absence from post-training, and couple embodiment changes with changes of task, scene and camera. Benchmarks such as LIBERO and its robustness extensions \citep{liu2023libero, zhou2025liberopro, fei2025liberoplus, wang2026liberox}, CALVIN \citep{mees2021calvin} and RoboTwin \citep{chen2025robotwin2} probe generalisation along task, language and layout axes on a fixed robot, and cross-embodiment suites report transfer to specific platforms. Our benchmark isolates embodiment shift in a controlled tabletop regime, separates the two zero-shot protocols and adds a diagnostic, the effect error on the target, that anticipates transfer before any task rollout.

\paragraph{Effect supervision.}
Point correspondence models \citep{karaev2024cotracker3} supply dense image motion for visual pretraining and for world models conditioned on predicted motion \citep{nvidia2025cosmospredict25}. Prior work in this line predicts the motion of scene points as an auxiliary signal on a single robot. We restrict the effect to the keypoints of the gripper, anchor it in the kinematics of the robot in hand, and study it as a cross-embodiment action in addition to its established role as a visual prediction target.
